\section{Introduction}

Simultaneous speech translation, or SST, generates target-language text while
a speaker is still talking, serving technical talks, lectures, and live media.
Speech-capable large language models, or LLMs, now translate long-form audio directly
\citep{xu2025qwen3omni,ouyang2025infinisst}. Rare method and dataset names, clinical
concepts, and named entities remain hard to recover from acoustics alone. Glossary retrieval helps by matching speech to source--target term pairs
and prompting the LLM with top matches before decoding.

Existing retrieval-augmented SST systems need manual session setup or lose
accuracy as terminology catalogs grow. Retrieval-Augmented Simultaneous
Speech Translation, or RASST \citep{luo2026rasst}, improves terminology on
out-of-domain talks, but fixes the glossary before translation. Listeners
rarely know a live talk's exact domain in advance. Finding, uploading, or
selecting a glossary before every session adds friction. Merging all domains
into one glossary avoids manual selection and works at small scale. With
fixed counts of retrieval candidates and prompt term pairs, irrelevant and
conflicting mappings crowd out relevant terms as the catalog grows. In our
evaluation, increasing the size of a merged glossary from 100,000 to one million
entries brings terminology accuracy down from 87.8\% to 84.4\%, as
Table~\ref{tab:quality} shows.

The translation stream itself supplies the evidence to choose a glossary. The
RealSI streams~\citep{cheng2024clasi} change domains every 3.4 to 5.2 minutes. Our
4-billion-parameter language-model router reads the last four translated chunks and
ranks a new talk's glossary first after a median of 15.6 seconds. A streaming translator can therefore infer the domain from its own
output without asking the listener to declare it.

We present AutoTerm-SST, a system demo that moves glossary selection from user
setup into the streaming loop. Operators register offline topic slices, each
holding roughly 1,000 term pairs and a prototype, a short description of the
topic. During a session, a small LLM scores every prototype against recent
partial translations, a sticky hidden Markov filter accumulates these scores
into a posterior over topics, and the most probable slices, at most four, form
a working set inside a fixed term budget. Retrieval and decoding stay
unchanged, so the per-chunk budget stays constant as the catalog grows. In
summary, our contributions are threefold:
\begin{itemize}
\item Zero-setup terminology-aware SST built on RASST, paired with a live
interface that displays the active slices and the retrieved term pairs beside
every partial translation.
\item Training-free language-model routing with a sticky topic filter whose
only free parameter is the switch rate. The correct slice stays searchable
in 89.7\% of chunks on ten-domain RealSI streams. Embedding-similarity
routing with tuned confirmation rules reaches 78.8\%.
\item Evidence from a 4.7-hour stream alternating natural language processing,
or NLP, and oncology talks: AutoTerm-SST is on par with a known-domain reference and beats a
one-million-entry merged glossary by 3.3 points of terminology accuracy and
2.31 BLEU.
\end{itemize}
